\documentclass[letterpaper,11pt,twoside]{article}
\usepackage{graphicx} % Required for inserting images
\usepackage[table,xcdraw,dvipsnames]{xcolor}
\usepackage{amsmath,amsfonts,amssymb,amsthm}
\usepackage{listings}
\usepackage{lipsum}
\usepackage{hyperref}
\usepackage{enumitem}

\usepackage{tikz}
\usepackage[siunitx, RPvoltages]{circuitikz}
\usetikzlibrary{3d}
\usepackage{comment}
\usepackage{caption,subcaption}
\usepackage{pgfplots}
\pgfplotsset{compat=newest} % or a newer version if available
\usepgfplotslibrary{groupplots}
\usetikzlibrary{pgfplots.groupplots}
\usetikzlibrary{shapes.geometric, arrows}
\tikzstyle{arrow} = [->,>=stealth,shorten >=2pt]
\newcommand{\re}[1]{\text{Re}\left(#1\right)}
\newcommand{\im}[1]{\text{Im}\left(#1\right)}
\newcommand{\bA}{\bm{A}}
\newcommand{\bB}{\bm{B}}
\newcommand{\bC}{\bm{C}}
\newcommand{\bE}{\bm{E}}
\newcommand{\bH}{\bm{H}}
\newcommand{\bD}{\bm{D}}
\newcommand{\bO}{\bm{0}}
\newcommand{\br}{\bm{r}}
\newcommand{\bk}{\bm{k}}
\newcommand{\bR}{\bm{R}}
\newcommand{\hA}{\hat{A}}
\newcommand{\hB}{\hat{B}}
\newcommand{\hC}{\hat{C}}
\newcommand{\ha}{\hat{a}}
\newcommand{\hb}{\hat{b}}
\newcommand{\hc}{\hat{c}}
\newcommand{\hn}{\hat{n}}
\newcommand{\hH}{\hat{H}}
\newcommand{\hN}{\hat{N}}
\newcommand{\hr}{\bm{\hat{r}}}
\newcommand{\hx}{\hat{x}}
\newcommand{\hX}{\hat{X}}
\newcommand{\hp}{\hat{p}}
\newcommand{\hS}{\hat{S}}
\newcommand{\hD}{\hat{D}}
\newcommand{\hy}{\hat{y}}
\newcommand{\hrho}{\hat{\rho}}
\newcommand{\hz}{\hat{z}}
\newcommand{\laplace}[1]{\mathscr{L}\left[#1\right]}
\newcommand{\ilaplace}[1]{\mathscr{L}^{-1}\left[#1\right]}
\newcommand{\fourier}[1]{\mathscr{F}\left[#1\right]}
\newcommand{\Tr}[1]{\text{Tr}\left[#1\right]}
\newcommand{\ifourier}[1]{\mathscr{F}^{-1}\left[#1\right]}
\newcommand{\ket}[1]{\left|#1\right\rangle}
\newcommand{\bra}[1]{\left\langle#1\right|}
\newcommand{\braket}[1]{\langle#1\rangle}
\newcommand{\D}{\mathcal{D}}
\usepackage{cancel}
\usepackage{bm}
\usepackage{fancyhdr}
\usepackage[utf8x]{inputenc}
\usepackage[T1]{fontenc}
\usepackage[margin=0.8in,top=1in,bottom=1in]{geometry}
%%%%%
\begin{filecontents*}{refs.bib}
@article{cohen1986quantum,
  title={Quantum mechanics, volume 1},
  author={Cohen-Tannoudji, Claude and Diu, Bernard and Laloe, Frank},
  journal={Quantum Mechanics},
  volume={1},
  pages={898},
  year={1986}
}
\end{filecontents*}
%
\newcommand{\institution}{University of Arizona}
\newcommand{\autor}{Nicolás Hernández Alegría}
\newcommand{\course}{OPTI 544 Quantum Optics}
\newcommand{\assignment}{Assignment 3}
%
\title{\textbf{\assignment}\\\course\\{\Large\institution}}
\author{\autor}
%\date{\today\\Total time: 12 hours}
%
\renewcommand{\sectionmark}[1]{\markright{#1}}
\fancypagestyle{mainstyle}{
    \fancyhf{} % Clear all header and footer fields
    \fancyfoot[C]{\thepage}
    \fancyhead[LE,RO]{\course} % Section name on odd pages
    \fancyhead[LO,RE]{\assignment}
    % Optional: Thin rules
    \renewcommand{\headrulewidth}{0pt} % Header rule
    \renewcommand{\footrulewidth}{0pt} % No footer rule
}
%
\begin{document}

\pagestyle{mainstyle}
\maketitle
%%
\section*{Exercise 1}
\begin{enumerate}[itemsep=0pt,topsep=0pt,parsep=0pt,label=\alph*)]
  \item sfasfasf
  \item sfasfasf
  \item asfasf
\end{enumerate}
%%
\section*{Exercise 2}
\begin{enumerate}[itemsep=0pt,topsep=0pt,parsep=0pt,label=\alph*)]
  \item First, the eigenequation of the number operator $\hn$ and the definition of the thermal density matrix are:
  \begin{align*}
    \hn\ket{n}=n\ket{n},\quad\text{and}\quad\hrho_{th}=\sum_nP_n\ket{n}\bra{n}.
  \end{align*}
  The expectation value of the number of photons is:
  \begin{align*}
    \braket{\hn}=\Tr{\hn\hrho_{th}}=\sum_n \braket{n|\hn\hrho_{th}|n}=\sum_{n,n'}\braket{n|\hn P_{n'}|n'}\braket{n'|n}=\sum_{n,n'}n'P_{n'}|\braket{n|n'}|^2=\sum_nnP_n.
  \end{align*}
  Now, we put the Boltzman distribution in the last summation:
  \begin{align*}
    \braket{\hn}=\sum_nn\left[\frac{e^{-(n+1/2)\hbar\omega/(k_BT)}}{\sum_{m=0}^\infty e^{-(m+1/2)\hbar\omega/(k_BT)}}\right]=\sum_nn\frac{e^{-n[\hbar\omega/(k_BT)]}}{\sum_{m=0}^\infty e^{-m[\hbar\omega/(k_BT)]}}=\frac{1}{\sum_{m=0}^\infty e^{-mx}}\sum_nne^{-nx},\;x=\frac{\hbar\omega}{k_BT}.
  \end{align*}
  We note that the summation in $n$ can be expressed in terms of a geometric function in the following way
  \begin{align*}
    \sum_{n=0}^\infty ne^{-nx}=-\frac{d}{dx}\left[\sum_{n=0}^\infty e^{-nx}\right]=-\frac{d}{dx}\left[\frac{1}{1-e^{-x}}\right]=\frac{e^{-x}}{(1-e^{-x})^2}.
  \end{align*}
  Then,
  \begin{align*}
    \braket{\hn}=\frac{1}{\sum_{m=0}^\infty e^{-mx}}\frac{e^{-x}}{(1-e^{-x})^2}=(1-e^{-x})\frac{e^{-x}}{(1-e^{-x})^2}=\frac{e^{-x}}{1-e^{-x}}=\frac{1}{e^{x}-1}=\frac{1}{e^{\frac{\hbar\omega}{k_BT}}-1}.
  \end{align*}
  which is the Planck radiation distribution.
  \item By evaluating the above result, and knowing that 
  \begin{align*}
    h\approx6.6261\cdot10^{-34}\;m^2kg/s,\quad\text{and}\quad k_B\approx1.3806\cdot10^{-23}\;m^2kgs^{-2}K^{-1}.
  \end{align*}
  The averague number of photon as a function of the absolute temperature is:
  \begin{align*}
    \braket{\hn}(T)=\frac{1}{e^{\dfrac{[6.6261\cdot10^{-34}/(2\pi)][(2\pi)(3\cdot10^8)/(500\cdot10^{-9})]}{(1.3806\cdot10^{-23})K}}-1}=\frac{1}{e^{28797/K}-1}.
  \end{align*}
  At $300K$ and $6000K$ we have 
  \begin{align*}
    \braket{\hn}(300)=\frac{1}{e^{28797/300}-1}=2\cdot10^{-42},\quad\text{and}\quad\braket{\hn}(6000)=\frac{1}{e^{28797/6000}-1}=8.3\cdot10^{-3}.
  \end{align*}
  The plot of the average number of photons along the temperature is given below.
\end{enumerate}

%\clearpage
%\nocite{*}
%\bibliographystyle{plain}   % or unsrt, alpha, apalike, etc.
%\bibliography{refs}

\end{document}
